\documentclass[aspectratio=169,8pt]{beamer}
\usetheme{Madrid}
\usepackage[utf8]{inputenc}
\usepackage[T1]{fontenc}
\usepackage{amsmath,amssymb}
\usepackage{booktabs}
\usepackage{enumitem}
\setlist[enumerate]{label=\Alph*.,leftmargin=*,itemsep=1pt,topsep=2pt}
\setlist[itemize]{leftmargin=*,itemsep=1pt,topsep=2pt}
\setbeamerfont{block body}{size=\tiny}
\setbeamerfont{block title}{size=\scriptsize}
\setbeamerfont{frametitle}{size=\footnotesize}
\setbeamerfont{normal text}{size=\tiny}
\setbeamertemplate{navigation symbols}{}
\setbeamertemplate{itemize item}{\tiny$\bullet$}
\setbeamertemplate{itemize subitem}{\tiny$\circ$}

\title[GARP RAI]{GARP Risk \& AI --- Q351--Q450}
\subtitle{Unofficial Exam Prep}
\author{Unofficial}
\date{\today}

\begin{document}
	
	\begin{frame}{Q540 --- Recall}
		\begin{block}{Question}
			Four fraud detection models are evaluated:
			
			\begin{center}
				\begin{tabular}{lccc}
					\hline
					Model & Precision & Recall & F1 \\
					\hline
					A & 0.80 & 0.22 & 0.34 \\
					B & 0.55 & 0.71 & 0.62 \\
					C & 0.77 & 0.41 & 0.54 \\
					D & 0.61 & 0.36 & 0.45 \\
					\hline
				\end{tabular}
			\end{center}
			
			Which model catches the largest percentage of actual fraud cases?
			\begin{enumerate}
				\item A
				\item B
				\item C
				\item D
			\end{enumerate}
		\end{block}
		
		\begin{exampleblock}{Answer: B}
			Model B
		\end{exampleblock}
		
		\begin{block}{Explanation}
			\textbf{Recall} (also called sensitivity or true positive rate) measures:
			\[
			\text{Recall} = \frac{TP}{TP+FN}
			\]
			``Of all the \emph{actual} fraud cases, what fraction did the model catch?''
			
			Comparing the recall column:
			\begin{itemize}
				\item A: 0.22
				\item \textbf{B: 0.71} \(\leftarrow\) highest
				\item C: 0.41
				\item D: 0.36
			\end{itemize}
			
			Model B has the highest recall (0.71), so it catches the largest percentage
			of actual fraud cases. Precision is irrelevant to this question---the question
			specifically asks about \emph{actual} fraud being caught, which is recall.
		\end{block}
	\end{frame}
	
	\begin{frame}{Q541 --- Precision}
		\begin{block}{Question}
			Which model produces the highest percentage of correct fraud alerts once it flags a transaction?
			
			\begin{center}
				\begin{tabular}{lcc}
					\hline
					Model & Precision & Recall\\
					\hline
					A & 0.91 & 0.18\\
					B & 0.62 & 0.70\\
					C & 0.71 & 0.52\\
					D & 0.68 & 0.44\\
					\hline
				\end{tabular}
			\end{center}
			
			\begin{enumerate}
				\item A
				\item B
				\item C
				\item D
			\end{enumerate}
		\end{block}
		
		\begin{exampleblock}{Answer: A}
			Model A
		\end{exampleblock}
		
		\begin{block}{Explanation}
			\textbf{Precision} answers:
			\[
			\text{Precision} = \frac{TP}{TP+FP}
			\]
			``Of all the transactions the model \emph{flagged} as fraud, how many were
			actually fraudulent?''
			
			The phrase ``once it flags a transaction'' is the key: it conditions on the
			\emph{prediction} being positive. That is exactly precision.
			
			Comparing the precision column:
			\begin{itemize}
				\item \textbf{A: 0.91} \(\leftarrow\) highest
				\item B: 0.62
				\item C: 0.71
				\item D: 0.68
			\end{itemize}
			
			Model A has the highest precision (0.91), meaning when it raises a fraud alert,
			it is correct 91\% of the time. Note that Model A has very low recall (0.18),
			so it misses most fraud---but the question only asks about correctness of
			alerts, not coverage.
		\end{block}
	\end{frame}
	
	\begin{frame}{Q542 --- Accuracy Trap}
		\begin{block}{Question}
			A loan portfolio contains 98\% good loans and 2\% defaults.
			
			A model predicts every loan as non-default.
			
			What is the model's approximate accuracy?
			\begin{enumerate}
				\item 2\%
				\item 50\%
				\item 98\%
				\item 100\%
			\end{enumerate}
		\end{block}
		
		\begin{exampleblock}{Answer: C}
			98\%
		\end{exampleblock}
		
		\begin{block}{Explanation}
			\textbf{Accuracy} is:
			\[
			\text{Accuracy} = \frac{TP + TN}{TP + TN + FP + FN}
			\]
			
			If the model predicts \emph{non-default} for every loan:
			\begin{itemize}
				\item It correctly classifies all 98\% good loans (true negatives).
				\item It incorrectly classifies all 2\% defaults as good (false negatives).
			\end{itemize}
			
			\[
			\text{Accuracy} = \frac{98\%}{100\%} = 98\%
			\]
			
			\textbf{The trap:} 98\% accuracy sounds excellent, but the model has
			\emph{zero} recall for defaults---it catches \emph{none} of the actual
			defaults. This is the classic \textbf{class imbalance problem}: a trivial
			majority-class predictor can achieve high accuracy while being completely
			useless for the minority (and usually more important) class.
			
			This is why accuracy alone is a poor metric for imbalanced problems, and
			why precision, recall, and F1 are preferred.
		\end{block}
	\end{frame}
	
	\begin{frame}{Q543 --- F1 Score}
		\begin{block}{Question}
			Which metric combines precision and recall into a single measure?
			
			\begin{enumerate}
				\item Accuracy
				\item FPR
				\item F1 Score
				\item AUC
			\end{enumerate}
		\end{block}
		
		\begin{exampleblock}{Answer: C}
			F1 Score
		\end{exampleblock}
		
		\begin{block}{Explanation}
			The \textbf{F1 score} is the \textbf{harmonic mean} of precision and recall:
			\[
			F1 = 2 \cdot \frac{\text{Precision} \cdot \text{Recall}}{\text{Precision} + \text{Recall}}
			\]
			
			\textbf{Why harmonic mean and not arithmetic mean?}
			The harmonic mean punishes extreme values. If a model has precision = 1.0
			and recall = 0.0, the arithmetic mean would be 0.5, but the harmonic mean
			is 0. So F1 is only high when \emph{both} precision and recall are
			reasonably high---it forces a balance.
			
			\textbf{Why the others are wrong:}
			\begin{itemize}
				\item \textbf{Accuracy:} measures overall correctness, not a
				precision/recall combination.
				\item \textbf{FPR:} false positive rate, a single error rate.
				\item \textbf{AUC:} area under the ROC curve, a threshold-independent
				ranking measure.
			\end{itemize}
		\end{block}
	\end{frame}
	
	\begin{frame}{Q544 --- F1 Ranking}
		\begin{block}{Question}
			Which model provides the best balance between precision and recall?
			
			\begin{center}
				\begin{tabular}{lccc}
					\hline
					Model & Precision & Recall & F1 \\
					\hline
					A & 0.90 & 0.15 & 0.26 \\
					B & 0.60 & 0.70 & 0.65 \\
					C & 0.75 & 0.30 & 0.43 \\
					D & 0.55 & 0.40 & 0.46 \\
					\hline
				\end{tabular}
			\end{center}
			
			\begin{enumerate}
				\item A
				\item B
				\item C
				\item D
			\end{enumerate}
		\end{block}
		
		\begin{exampleblock}{Answer: B}
			Model B
		\end{exampleblock}
		
		\begin{block}{Explanation}
			``Best balance between precision and recall'' means the highest \textbf{F1 score},
			since F1 is the harmonic mean of the two.
			
			Comparing the F1 column:
			\begin{itemize}
				\item A: 0.26 (high precision 0.90 but terrible recall 0.15)
				\item \textbf{B: 0.65} \(\leftarrow\) highest
				\item C: 0.43
				\item D: 0.46
			\end{itemize}
			
			Model B has the highest F1 (0.65), meaning it achieves the best trade-off
			between catching positives (recall 0.70) and being correct when it predicts
			positive (precision 0.60).
			
			\textbf{Note:} Model A has the highest precision (0.90) but the \emph{worst}
			F1, because its recall is only 0.15. This illustrates why you should not
			judge a model by a single metric---high precision alone does not mean a
			good balance.
		\end{block}
	\end{frame}
	
	\begin{frame}{Q545 --- ROC Curve}
		\begin{block}{Question}
			An ROC curve plots:
			
			\begin{enumerate}
				\item Precision against Recall
				\item TPR against FPR
				\item Accuracy against Threshold
				\item Recall against Accuracy
			\end{enumerate}
		\end{block}
		
		\begin{exampleblock}{Answer: B}
			TPR against FPR
		\end{exampleblock}
		
		\begin{block}{Explanation}
			An \textbf{ROC curve} (Receiver Operating Characteristic) plots:
			\[
			\text{True Positive Rate (TPR)} \quad \text{vs.} \quad \text{False Positive Rate (FPR)}
			\]
			across all possible classification thresholds.
			
			Where:
			\[
			\text{TPR} = \frac{TP}{TP+FN} = \text{Recall / Sensitivity}
			\]
			\[
			\text{FPR} = \frac{FP}{FP+TN} = 1 - \text{Specificity}
			\]
			
			\textbf{How to read it:}
			\begin{itemize}
				\item The top-left corner (TPR = 1, FPR = 0) is perfect classification.
				\item The diagonal line (TPR = FPR) represents random guessing.
				\item A curve closer to the top-left corner is a better model.
			\end{itemize}
			
			\textbf{Why the others are wrong:}
			\begin{itemize}
				\item \textbf{A:} Precision vs. Recall is the \emph{Precision-Recall curve},
				a different (and often more informative) plot for imbalanced data.
				\item \textbf{C:} Accuracy vs. Threshold is a threshold-sweep plot, not ROC.
				\item \textbf{D:} Recall vs. Accuracy is not a standard curve.
			\end{itemize}
		\end{block}
	\end{frame}
	
	\begin{frame}{Q546 --- AUC Interpretation}
		\begin{block}{Question}
			Two models have AUC values:
			
			\[
			AUC_A = 0.94
			\]
			
			\[
			AUC_B = 0.72
			\]
			
			Which statement is most accurate?
			
			\begin{enumerate}
				\item Model B separates classes better.
				\item Model A separates classes better.
				\item Both models perform identically.
				\item AUC measures precision.
			\end{enumerate}
		\end{block}
		
		\begin{exampleblock}{Answer: B}
			Model A separates classes better.
		\end{exampleblock}
		
		\begin{block}{Explanation}
			\textbf{AUC} (Area Under the ROC Curve) measures the model's ability to
			\emph{discriminate} between positive and negative classes.
			
			\textbf{Interpretation:} AUC is the probability that a randomly chosen
			positive example is ranked higher (scored more positively) than a randomly
			chosen negative example.
			
			\textbf{Range and meaning:}
			\begin{itemize}
				\item AUC = 0.5: random guessing (no discrimination)
				\item AUC = 1.0: perfect separation
				\item AUC = 0.94 (Model A): excellent discrimination
				\item AUC = 0.72 (Model B): acceptable discrimination
			\end{itemize}
			
			Since \(AUC_A = 0.94 > AUC_B = 0.72\), \textbf{Model A separates the classes
				better}.
			
			\textbf{Why the others are wrong:}
			\begin{itemize}
				\item \textbf{A:} Model B has the \emph{lower} AUC, so it separates worse.
				\item \textbf{C:} The AUCs are clearly different, so they do not perform identically.
				\item \textbf{D:} AUC is not precision. Precision is a single-threshold
				metric; AUC is threshold-independent.
			\end{itemize}
		\end{block}
	\end{frame}
	
	\begin{frame}{Q547 --- Threshold Effects}
		\begin{block}{Question}
			A bank lowers its approval threshold for identifying risky loans.
			
			What is the most likely outcome?
			
			\begin{enumerate}
				\item Lower recall and lower FPR
				\item Higher recall and higher FPR
				\item Higher precision and lower recall
				\item Lower precision and lower recall
			\end{enumerate}
		\end{block}
		
		\begin{exampleblock}{Answer: B}
			Higher recall and higher FPR
		\end{exampleblock}
		
		\begin{block}{Explanation}
			Lowering the threshold means the model becomes \textbf{more eager to predict
				the positive class} (``risky loan''). It flags more loans as risky.
			
			\textbf{Effect on recall (TPR):}
			\begin{itemize}
				\item More loans are flagged positive.
				\item More of the truly risky loans get caught.
				\item \(\Rightarrow\) Recall \textbf{increases}.
			\end{itemize}
			
			\textbf{Effect on FPR:}
			\begin{itemize}
				\item More loans are flagged positive overall.
				\item More \emph{safe} loans get incorrectly flagged as risky.
				\item \(\Rightarrow\) FPR \textbf{increases}.
			\end{itemize}
			
			This is the fundamental \textbf{threshold trade-off}: lowering the threshold
			increases both TPR and FPR. You cannot increase recall without also
			increasing false positives (unless the model is perfect).
			
			\textbf{Why the others are wrong:}
			\begin{itemize}
				\item \textbf{A:} Recall does not decrease when you lower the threshold.
				\item \textbf{C:} Precision typically \emph{decreases} (more false positives),
				and recall \emph{increases}, not the reverse.
				\item \textbf{D:} Recall increases, not decreases.
			\end{itemize}
		\end{block}
	\end{frame}
	
	\begin{frame}{Q548 --- False Positive Rate}
		\begin{block}{Question}
			A model incorrectly flags 200 legitimate transactions out of 5,000 legitimate transactions.
			
			The FPR is closest to:
			
			\begin{enumerate}
				\item 1\%
				\item 4\%
				\item 20\%
				\item 40\%
			\end{enumerate}
		\end{block}
		
		\begin{exampleblock}{Answer: B}
			4\%
		\end{exampleblock}
		
		\begin{block}{Explanation}
			\textbf{False Positive Rate (FPR)} is:
			\[
			FPR = \frac{FP}{FP + TN}
			\]
			
			Here:
			\begin{itemize}
				\item \(FP = 200\) (legitimate transactions incorrectly flagged)
				\item \(TN = 5000 - 200 = 4800\) (legitimate transactions correctly not flagged)
				\item Total legitimate transactions (actual negatives) \(= 5000\)
			\end{itemize}
			
			\[
			FPR = \frac{200}{5000} = 0.04 = 4\%
			\]
			
			\textbf{Why the others are wrong:}
			\begin{itemize}
				\item \textbf{A (1\%):} Would require only 50 false positives.
				\item \textbf{C (20\%):} Would require 1,000 false positives.
				\item \textbf{D (40\%):} Would require 2,000 false positives.
			\end{itemize}
			
			\textbf{Note:} The denominator is the total number of \emph{actual negatives}
			(all legitimate transactions), not the number of predicted positives. This is
			a common point of confusion---do not use \(FP/(FP+TP)\) (that would be
			\(1 - \text{precision}\)).
		\end{block}
	\end{frame}
	
	\begin{frame}{Q549 --- Metric Matching}
		\begin{block}{Question}
			Which pairing is correct?
			
			\begin{enumerate}
				\item Precision = Of actual positives, how many were found?
				\item Recall = Of predicted positives, how many were correct?
				\item Precision = Of predicted positives, how many were correct?
				\item FPR = Of actual positives, how many were missed?
			\end{enumerate}
		\end{block}
		
		\begin{exampleblock}{Answer: C}
			Precision = Of predicted positives, how many were correct?
		\end{exampleblock}
		
		\begin{block}{Explanation}
			The correct definitions:
			\[
			\text{Precision} = \frac{TP}{TP + FP} \quad \text{(of \emph{predicted} positives, how many are correct?)}
			\]
			\[
			\text{Recall} = \frac{TP}{TP + FN} \quad \text{(of \emph{actual} positives, how many are found?)}
			\]
			
			\textbf{Why the others are wrong:}
			\begin{itemize}
				\item \textbf{A:} ``Of actual positives, how many were found?'' is the
				definition of \textbf{recall}, not precision.
				\item \textbf{B:} ``Of predicted positives, how many were correct?'' is the
				definition of \textbf{precision}, not recall. (The pairing is swapped.)
				\item \textbf{D:} ``Of actual positives, how many were missed?'' is the
				\textbf{false negative rate} \((FN/(TP+FN) = 1 - \text{recall})\),
				not FPR. FPR is ``of actual \emph{negatives}, how many were
				incorrectly flagged?''
			\end{itemize}
			
			\textbf{Memory aid:}
			\begin{center}
				\begin{tabular}{ll}
					\hline
					\textbf{Metric} & \textbf{Conditions on} \\
					\hline
					Precision & Predicted positives (\(\hat{Y}=1\)) \\
					Recall & Actual positives (\(Y=1\)) \\
					FPR & Actual negatives (\(Y=0\)) \\
					\hline
				\end{tabular}
			\end{center}
		\end{block}
	\end{frame}
	
	\begin{frame}{Q541 --- Precision}
		\begin{block}{Question}
			Which model produces the highest percentage of correct fraud alerts once it flags a transaction?
			
			\begin{center}
				\begin{tabular}{lcc}
					\hline
					Model & Precision & Recall\\
					\hline
					A & 0.91 & 0.18\\
					B & 0.62 & 0.70\\
					C & 0.71 & 0.52\\
					D & 0.68 & 0.44\\
					\hline
				\end{tabular}
			\end{center}
			
			\begin{enumerate}
				\item A
				\item B
				\item C
				\item D
			\end{enumerate}
		\end{block}
		
		\begin{exampleblock}{Answer: A}
			Model A
		\end{exampleblock}
		
		\begin{block}{Explanation}
			\textbf{Precision} answers:
			\[
			\text{Precision} = \frac{TP}{TP+FP}
			\]
			``Of all the transactions the model \emph{flagged} as fraud, how many were
			actually fraudulent?''
			
			The phrase ``once it flags a transaction'' is the key: it conditions on the
			\emph{prediction} being positive. That is exactly precision.
			
			Comparing the precision column:
			\begin{itemize}
				\item \textbf{A: 0.91} \(\leftarrow\) highest
				\item B: 0.62
				\item C: 0.71
				\item D: 0.68
			\end{itemize}
			
			Model A has the highest precision (0.91), meaning when it raises a fraud alert,
			it is correct 91\% of the time. Note that Model A has very low recall (0.18),
			so it misses most fraud---but the question only asks about correctness of
			alerts, not coverage.
		\end{block}
		
		\begin{block}{Why A is correct and B is wrong}
			\textbf{A is correct.} The question asks about the percentage of
			\emph{correct fraud alerts once a transaction is flagged}. This conditions
			on the model having \emph{predicted} positive, which is exactly the
			definition of \textbf{precision}. Model A has the highest precision
			(0.91), so it is the correct answer.
			
			\textbf{B is incorrect.} Model B has the highest recall (0.70)---meaning
			it catches the most actual fraud cases (sensitivity / true positive rate).
			But recall \emph{ignores false positives}. A model can have high recall
			by flagging almost everything as fraud, which would produce many
			incorrect alerts. Precision, not recall, measures the correctness of
			the alerts that were actually raised.
			
			\textbf{Key distinction:}
			\begin{itemize}
				\item \textbf{Precision} = of predicted positives, how many were
				correct? (used here --- ``once it flags'')
				\item \textbf{Recall / Sensitivity} = of actual positives, how many
				were found? (used when missing a true case is the concern)
			\end{itemize}
			
			In a scenario where \emph{missing a true fraud case is fatal}, you would
			prioritize \textbf{recall} (sensitivity) and choose Model B. But in a
			scenario where \emph{every alert must be trustworthy} (e.g., to avoid
			annoying customers with false alarms), you prioritize \textbf{precision}
			and choose Model A.
		\end{block}
	\end{frame}
	
\end{document}